\section{Введение}
\label{sec:Chapter0}

Сердечно-сосудистые заболевания занимают первое место среди причин смертности
в~мире: по~данным Всемирной организации здравоохранения, ежегодно от~них
погибает более~17~миллионов человек~\cite{who2021}.
Значительную долю среди них составляют нарушения сердечного ритма~--- аритмии,
которые нередко протекают бессимптомно, но~при~отсутствии своевременной диагностики
могут приводить к~инфаркту миокарда, инсульту или~внезапной остановке сердца.
Главным инструментом выявления аритмий является электрокардиограмма~(ЭКГ)~---
запись электрической активности сердца во~времени.

При~суточном холтеровском мониторировании за~24~часа регистрируется от~80\,000
до~120\,000 сердечных сокращений.
Ручной разбор такой записи квалифицированным кардиологом требует многих часов работы
и~неизбежно сопровождается усталостными ошибками.
Именно поэтому автоматическая классификация ЭКГ-сигналов средствами машинного
обучения представляет большой практический интерес: она позволяет ускорить
диагностику и~снизить риск пропуска опасных состояний.

Особую сложность представляет \textit{нестационарность} ЭКГ-сигнала~---
его амплитуда, форма и~ритм меняются со~временем под~влиянием физической нагрузки,
дыхания и~эмоционального состояния пациента.
Классические методы спектрального анализа Фурье, рассчитанные на~стационарные сигналы,
не~позволяют в~полной мере описать такую изменчивость.
Более гибким инструментом в~данном случае является вейвлет-преобразование,
которое одновременно учитывает частотный состав сигнала и~момент его
возникновения во~времени~\cite{zakharova2018}.

За~последние годы глубокое обучение существенно продвинулось в~задаче
распознавания аритмий: сверточные и~остаточные нейронные сети достигают точности,
сопоставимой с~экспертом-кардиологом~\cite{hannun2019}.
Тем не менее традиционные методы~--- в~первую очередь случайный лес~---
по-прежнему остаются востребованными благодаря высокой интерпретируемости
и~устойчивости к~помехам, что~особенно ценно в~клинической практике,
где надежность и~объяснимость модели не~менее важны, чем метрики точности.

Настоящая работа посвящена сравнению трех подходов к~автоматической классификации
кардиоциклов~(нормы и~четырех типов аритмий) на~базе данных MIT-BIH Arrhythmia
Database: случайного леса, одномерной сверточной нейронной сети и~остаточной
нейронной сети.
Помимо классификации на~чистых данных, исследуется устойчивость моделей
к~четырем видам реалистичных помех~--- гауссову шуму, дрейфу изолинии,
сетевой помехе~50~Гц и~мышечному шуму~--- а~также проводится
частотно-временной анализ сигналов с~помощью дискретного вейвлет-преобразования.

Работа состоит из~пяти разделов.
В~разделе~\ref{sec:Chapter1} формализуется постановка задачи и~описываются
данные и~конвейер предобработки.
Раздел~\ref{sec:Chapter2} посвящен теоретическому описанию моделей и~методам
вейвлет-анализа.
В~разделе~\ref{sec:Chapter3} приводятся результаты классификации на~чистых данных.
Раздел~\ref{sec:Chapter4} содержит анализ устойчивости к~шуму и~результаты
кросс-валидации.
В~заключении~(\ref{sec:Chapter5}) формулируются выводы и~направления
дальнейших исследований.
